\BeleskeID{08-s048}
% element: 08-s048:e02
% element: 08-s048:notes
Iz uslova za zbir konstanti i početni nagib dobija se $K_1=-(U-U_0)$ i $K_2=U$. Dva zapisa izlaznog napona razlikuju se samo algebarskim rasporedom članova: prvi ističe približavanje krajnjoj vrednosti $U$, a drugi ostvareni deo ukupne promene $U-U_0$.

U početku je eksponencijalni faktor jedan i izlaz zadržava početni napon. Tokom vremena eksponencijalni član iščezava. Ovo računanje važi u intervalu u kome je nova pobuda konstantna.
